\begin{frontmatter}

\title{Frequent Losers as Cover Bidders in Public Procurement\texorpdfstring{\footnote{This version: April 2026.}}{}}

\author[insper]{Darcio Genicolo-Martins\corref{cor1}}
\ead{darciogm1@insper.edu.br}
\author[insper]{Paulo Furquim de Azevedo}
\affiliation[insper]{organization={INSPER},
  city={S\~ao Paulo}, country={Brazil}}
\cortext[cor1]{Corresponding author.}

\begin{abstract}
Bid-rigging cartels must simulate competition to avoid detection,
often deploying firms that submit deliberately losing bids. We
propose a firm-level screening marker based on this footprint:
\emph{frequent losers} (FL)---firms that never win yet participate
in abnormally many tenders. Using 4.5 million tender-items from
Brazilian public procurement (2009--2019), we document a robust
conditional association between FL presence and higher prices:
temporal cross-fitting---which breaks any mechanical link between
the classification and the outcome---yields an out-of-sample
estimate of 3.6\%; full-sample OLS with high-dimensional fixed
effects yields 6.4\%; matching estimators bracket these figures
(5.5--7.7\%). External validation
shows that 7.1\% of FL firms co-participate with
competition-authority-convicted cartelists---3.5 times the expected
rate ($p < 0.001$). Multiple diagnostics---bid-coordination tests,
co-bidding network analysis, and structural estimation---provide
suggestive evidence that cover bidding drives the association. The
screen requires only participation and outcome records, making it
operable where existing bid-level forensic tools cannot reach.
\end{abstract}

\end{frontmatter}

\noindent\textbf{Keywords:} bid rigging, cover bidding, public procurement, cartel screening, frequent losers

\noindent\textbf{JEL No.:} D44, D73, H57, K21, L41

\bigskip
